\section{Conclusiones y limitaciones}
\label{sec:conclusiones}

\paragraph{Diferencias entre modelos.}
En la partición de prueba de 4x4 y 6x6, el 9B supera al 2B en casi todas las
celdas, y la ventaja se concentra donde el 2B empieza a fallar: en 4x4 con 12
casillas, \passuno\ de 68,0\,\% frente a 14,0\,\% sin ejemplos, y en 6x6 con 8
casillas, de 86,0\,\% frente a 21,0\,\% (sección~\ref{sec:prueba}). Promediando
los seis niveles, \passuno\ es 31,0\,\% para el 2B y 59,7\,\% para el 9B sin
ejemplos.

\paragraph{Efecto de los ejemplos y del límite de tokens.}
Dos \emph{shots} mejoran el 2B (de 31,0\,\% a 39,5\,\% de \passuno\ promedio) y
reducen los tokens generados en los niveles fáciles, pero no abren los niveles sin
aciertos. El fallo dominante en toda la evaluación es el truncamiento: el 50,1\,\%
de las 7.680 respuestas alcanzó el límite de 2.304 tokens, y casi todas las
respuestas de 6x6 con 16 o 24 casillas (sección~\ref{sec:cualitativo}). Parte de
los fallos refleja el presupuesto de tokens y no la capacidad.

\paragraph{Robustez del verificador.}
El verificador acepta cualquier solución válida que conserve las pistas y
rechaza los casos probados, y los 7.680 \emph{rewards} registrados coinciden con
los recalculados (sección~\ref{sec:robustez}). Las reglas del parser sobre el
razonamiento delimitado no tienen pruebas propias.

\paragraph{Calibración del entorno.}
Para el 2B, la zona con señal de aprendizaje son 4x4 con 8 y 12 casillas y 6x6 con
8 casillas, donde casi todas las tareas tienen resultados mixtos; 4x4 con 4
casillas está casi saturado con ejemplos, y 6x6 con 16 y 24 casillas no tiene
aciertos para el 2B sin ejemplos (sección~\ref{sec:calibracion}).

\paragraph{Limitaciones.}
Se usó un solo \emph{seed}, sin intervalos de confianza, con 20 tareas por nivel.
Los tableros de 9x9 de la partición de prueba no se evaluaron, y la partición de
validación solo se usó en pilotos del 9B con muy pocas tareas, por lo que no hay
una línea base sobre ella. El límite de 2.304 tokens trunca la mayoría de las
respuestas en los niveles difíciles. Las \emph{traces} de entrenamiento
(sección~\ref{sec:traces}) se generaron, pero no se entrenó con ellas.

\paragraph{Cambios prioritarios antes del entrenamiento.}
Evaluar el 2B con un límite de tokens mayor en 6x6 para separar el truncamiento
del error de razonamiento, ejecutar los tableros de 9x9 de la prueba, repetir con
más \emph{seeds} y entrenar al 2B con las \emph{traces} ya generadas (SFT y DPO).
